% SET ARTICLE STYLE
\documentclass[paper=a4,fontsize=11pt,twocolumn]{article}


% PACKAGES
\usepackage[english]{babel} % English language/hyphenation
%\usepackage{fontspec} % To use Avenir
\usepackage{geometry} % For paper size and layout
\usepackage{color}
%\usepackage{xcolor}
\usepackage[svgnames]{xcolor} % Enabling colors by their 'svgnames'
\usepackage{fancyhdr} % For header and footer with odd and even page params
\usepackage{lipsum}  % For generating sample text
\usepackage{ragged2e} % For restating fully justified environment
\usepackage{textcomp} % ??? Can't remember...
\usepackage{titlesec} % To edit title formats
\usepackage{lastpage} % Ffor Page X of Y
\usepackage{booktabs} % For nicer tables
\usepackage{multicol} % More things for tables (and arranging text)
\usepackage{graphicx} % For easier figure wrangling
\usepackage{lineno} % Line numbers for review purposes
\usepackage{microtype} % For uneven columns
\usepackage{vwcol} % For uneven columns
\usepackage[T1]{fontenc} % For fonts
\usepackage{lmodern} % {tgheros lmodern paratype} For fonts
\usepackage{apacite} % For citations
\renewcommand{\familydefault}{\sfdefault}

% DEFINITIONS
\linenumbers
 \geometry{
	 a4paper,
	 top=20mm,
	 right=15mm,
	 left=15mm,
 	 bottom=20mm,
 	 }
\definecolor{brand-primary}{HTML}{1A2530}
\definecolor{brand-secondary}{HTML}{2C5A8F}
\definecolor{brand-logo-dark}{HTML}{58778C}
\definecolor{brand-logo-light}{HTML}{BCD3DB}
\definecolor{brand-midlight}{HTML}{C2C5C3}

\titleformat*{\section}{\Large\normalfont}
\titleformat*{\subsection}{\large\normalfont}
\titleformat*{\subsubsection}{\large\normalfont}
\titleformat*{\paragraph}{\normalfont}


% HEADERS AND FOOTERS 
\pagestyle{fancy}
\fancyhf{}
\fancyhead[RE,LO]{ARTICLES}

\fancyhead[LE,RO]{\footnotesize\textcolor{brand-secondary}{Running header of this paper}}
\fancyfoot[LE,LO]{\footnotesize{\textcopyright \textcolor{brand-secondary}{Journal of the Student Association of Medical Research}\\
\hspace{4mm}\footnotesize2017 Vol 1 Issue 1 \emph{doi:xx.xxxx/xxx}}}
\fancyfoot[RE,RO]{\footnotesize Page \thepage\ of \pageref{LastPage}}
\renewcommand{\headrulewidth}{2pt}
\renewcommand{\footrulewidth}{1pt}


% MISC
\graphicspath{{figures/}}
\hbadness=100000 % Turn off those dreadful overfull and underful hboxes
\hfuzz=1000pt
\vbadness=100000
\vfuzz=1000pt
\bibliographystyle{apacite}


% MACROS
\newcommand{\horizontalstrip}[1]{\par\noindent\colorbox{shadecolor}
{\parbox{\dimexpr\textwidth-2\fboxsep\relax}{#1}}}



% START OF DOCUMENT
\begin{document}
\pagestyle{fancy}

\twocolumn[

	\begin{@twocolumnfalse}	

		% Title 
		\huge\textcolor{brand-secondary}{Title of this paper}\\	
		
		% Author(s)
		\Large{First Author$^{*}$, Second Author$^{**}$ and Third Author$^{***}$}	

		% Affiliation(s)
		\setlength{\columnseprule}{1pt}
		\def\columnseprulecolor{\color{brand-secondary}}
		\begin{multicols}{2}
			\small{$^{*\alpha}$affiliation\smallskip\\
				$^{**}$affiliation\smallskip\\
				$^{***}$affiliation}\bigskip\smallskip\\
			\vfill
			\small{$^{\alpha}$Corresponding author: email@domain.ac.uk}
			\columnbreak

		% Dates	
			\hangindent=0.8cm 
			\small{
				\hspace{8mm}Received: 16$^{th}$ January 2017\smallskip\\
				Revised: 31$^{st}$ January 2017, 10$^{th}$ February 2017\smallskip\\
				Accepted: 11$^{th}$ February 2017\smallskip\\
				Available online: 12$^{th}$ February 2017\bigskip\\
				\vfill
				\hspace{8mm}Keywords: \colorbox{brand-midlight}{\strut keyword 1} \colorbox{brand-midlight}{\strut keyword 2} \colorbox{brand-midlight}{\strut keyword 3}
				}
		\end{multicols}
		\bigskip		

		% Abstract
		\begin{minipage}{0.95\textwidth}
		\begin{centering}
		\setlength{\fboxsep}{10pt}
		\fcolorbox{brand-primary}{brand-logo-light}{\parbox{\textwidth}{
			\Large{Abstract}\bigskip\\
			\normalsize
			\textbf{Introduction}
				Lorem ipsum dolor sit amet, consectetuer adipiscing elit. Aenean commodo ligula eget dolor. Aenean massa. Cum sociis natoque penatibus et magnis dis parturient montes, nascetur ridiculus mus. Donec quam felis, ultricies nec, pellentesque eu, pretium quis, sem.\\
			\textbf{Method}
				Lorem ipsum dolor sit amet, consectetuer adipiscing elit. Aenean commodo ligula eget dolor. Aenean massa. Cum sociis natoque penatibus et magnis dis parturient montes, nascetur ridiculus mus. Donec quam felis, ultricies nec, pellentesque eu, pretium quis, sem.\\
			\textbf{Results}
				Lorem ipsum dolor sit amet, consectetuer adipiscing elit. Aenean commodo ligula eget dolor. Aenean massa. Cum sociis natoque penatibus et magnis dis parturient montes, nascetur ridiculus mus. Donec quam felis, ultricies nec, pellentesque eu, pretium quis, sem.\\
			\textbf{Conclusions}
				Lorem ipsum dolor sit amet, consectetuer adipiscing elit. Aenean commodo ligula eget dolor. Aenean massa. Cum sociis natoque penatibus et magnis dis parturient montes, nascetur ridiculus mus. Donec quam felis, ultricies nec, pellentesque eu, pretium quis, sem.\\
			}
		} % End fbox
		\end{centering}
		\end{minipage} % END ABSTRACT

		
	\end{@twocolumnfalse}
\vspace{8mm}
] % End @twocolumnfalse


% Article text
\section{Introduction}
\lipsum[1-7]
\cite{Buss2008}

\section{Methods}
\subsection{This is a subheading}
Figure fig:one-column shows a one-column figure.

\begin{figure}
\centering
\includegraphics[width=\columnwidth]{logo-wide-jsamr.pdf}
\caption{Figure legend goes here if required. This is a single-column figure. All figures must be vector graphics if possible, or at least 600dpi.}
\label{fig:one-column}
\end{figure}

\lipsum[8-9]

\subsection{This is a subheading}
\lipsum[10-11]
\cite {Burns1981}

\section{Results}
\subsection{This is a subheading}
\subsubsection{This is a sub-subheading}

Table \ref{tab:one-column} shows an example of a simple single column table.

\begin{table}
\begin{centering}
\caption{A single-column table}\label{tab:one-column}
	\begin{tabular}[\columnwidth]{ccl}\toprule
		Experiment &  Variable 1 &  Variable 2 \\
		\midrule
		1  & X & A \\
		2  & X & A, B \\
		3  & X & C \\
		4  & Y & A \\
		5  & Y & A, B \\
		6  & Y & C \\
		\bottomrule
	\end{tabular}\par
	\medskip
\textit{Notes for this table if required. All values should be rounded to 2 dp}
\end{centering}
\end{table}

Table \ref{tab:two-column} shows an example of a table that requires two columns.

\begin{table*}
\begin{centering}
\caption{A two-column table}\label{tab:two-column}%
	\begin{tabular}[\columnwidth]{ccccccccc}\toprule
	Condition & V1 & V2 & V3 & V4 & V5 & V6 & V7 & V8 \\
	\midrule
	1 & 0.97 & --0.38 & 0.02 & 0.53 & --0.79 & --0.21 & 0.36 & --0.46 \\
	2 & 0.97 & --0.83 & 0.12 & --0.50 & 0.70 & --0.80 & --0.14 & --0.29 \\
	3 & 0.57 & --0.90 & 0.52 & --0.47 & --0.43 & --0.01 & --0.20 & --0.15 \\
	4 & 0.89 & --0.90 & 0.39 & --0.11 & --0.31 & --0.93 & 0.62 & --0.46 \\
	5 & 0.73 & --0.08 & 0.48 & --0.18 & 0.10 & 0.01 & --0.29 & --0.15 \\
	6 & --0.98 & --0.68 & --0.35 & 0.01 & --0.58 & --0.02 & --0.31 & 0.39 \\
	\cmidrule{2-9}
	Mean & 0.00 & --0.00 & 0.00 & --0.00 & 0.00 & --0.00 & --0.00 & --0.00 \\
	\bottomrule
	\end{tabular}\par
	\medskip
\textit{Notes for this table if required. All values should be rounded to 2 dp}
\end{centering}
\end{table*}

\lipsum[12]
\subsubsection{This is a sub-subheading}
\lipsum[13]

\subsection{This is a subheading}
\subsubsection{This is a sub-subheading}
\lipsum[14]
\subsubsection{This is a sub-subheading}
\lipsum[15]

\subsection{This is a subheading}
\lipsum[16-18]

\section{Discussion}
\lipsum[19-23]

\section{Conclusions}
\lipsum[24-32]

\bibliography{sample-article}

\end{document}

